\section{Вычислительный эксперимент в математическом моделировании}
\label{sec:computational_experiment}

\textbf{Вычислительный эксперимент} --- исследование математической модели
объекта или процесса посредством численных методов и компьютерной
программы. В отличие от натурного эксперимента непосредственно исследуется
не реальный объект, а его математическая модель.

Удобно выделять триаду
\[
\boxed{
\text{математическая модель}
\longrightarrow
\text{вычислительный алгоритм}
\longrightarrow
\text{программа}
}.
\]
Каждый уровень вносит собственные предположения и погрешности, поэтому
получение численного результата само по себе ещё не доказывает его
физическую достоверность.

\subsection{Математическая и вычислительная модели}

Пусть математическая модель записана абстрактно как
\[
\mathcal M(u,p)=0,
\]
где $u$ --- искомое состояние, $p$ --- параметры модели. После
дискретизации возникает конечномерная задача
\[
\mathcal M_h(u_h,p)=0,
\]
где $h$ условно обозначает параметры дискретизации: шаг сетки, шаг по
времени, число конечных элементов и т.п.

Вычислительный эксперимент состоит не в одном расчёте, а в
\textbf{серии контролируемых расчётов}, с помощью которых изучают
\begin{itemize}
    \item влияние параметров $p$ на выходные характеристики;
    \item переходные и стационарные режимы;
    \item критические значения параметров и потерю устойчивости;
    \item чувствительность и неопределённость результата;
    \item оптимальные режимы;
    \item область применимости математической модели.
\end{itemize}

\subsection{Основные этапы вычислительного эксперимента}

Типичная последовательность имеет вид:
\begin{enumerate}
    \item \textbf{Постановка цели исследования.} Определяются наблюдаемые
    величины и требуемая точность.
    \item \textbf{Построение математической модели.} Формулируются
    уравнения, определяющие соотношения, начальные и граничные условия,
    ограничения и допущения.
    \item \textbf{Выбор численного метода.} Непрерывная задача заменяется
    дискретной, анализируются аппроксимация, устойчивость и сходимость.
    \item \textbf{Программная реализация и верификация.} Проверяется, что
    программа действительно решает сформулированную математическую задачу;
    используются, в частности, аналитические решения, эталонные тесты и
    исследования сеточной/временной сходимости.
    \item \textbf{Планирование серии расчётов.} Выбираются диапазоны
    параметров, расчётные точки и критерии сравнения.
    \item \textbf{Проведение расчётов и постобработка.} Строятся графики,
    поля, интегральные характеристики, зависимости от параметров.
    \item \textbf{Валидация и интерпретация.} Результаты сопоставляются с
    независимыми экспериментальными, натурными или наблюдательными данными о
    реальном объекте и формулируются выводы о применимости модели.
\end{enumerate}

\subsection{Планирование серии расчётов}

Пусть выходная характеристика модели имеет вид
\[
y=F(p_1,\ldots,p_m).
\]
Простейший способ исследования --- изменять один параметр при фиксированных
остальных. Однако при наличии взаимодействий параметров такой подход может
быть недостаточен. Тогда применяют идеи \textbf{планирования эксперимента}:
выбирают набор точек в пространстве параметров, позволяющий с минимальным
числом расчётов оценить основные эффекты и построить аппроксимирующую
зависимость
\[
y\approx\widehat F(p_1,\ldots,p_m).
\]
Возможны полные и дробные факторные планы, планы поверхности отклика,
латинские гиперкубы и другие схемы. Для экзамена важно не перечисление
всех планов, а понимание идеи: расчётные точки выбираются не случайно, а
так, чтобы получить максимум информации о модели.

\subsection{Пример: теплопроводность стержня}

Рассмотрим
\[
\frac{\partial T}{\partial t}
=\alpha\frac{\partial^2T}{\partial x^2},
\qquad 0<x<L,
\]
с начальными и граничными условиями, например
\[
T(x,0)=T_0(x),\qquad
T(0,t)=T_1,
\qquad
\frac{\partial T}{\partial x}(L,t)=0.
\]
После введения сетки $x_i=ih$, $t^n=n\tau$ можно использовать, например,
разностную аппроксимацию
\[
\frac{T_i^{n+1}-T_i^n}{\tau}
=\alpha\frac{T_{i-1}^n-2T_i^n+T_{i+1}^n}{h^2}.
\]
Затем проводится не один расчёт, а серия: изменяют $\alpha$, начальную
температуру, длину стержня и другие параметры. Одновременно обязательно
проверяют зависимость результата от $h$ и $\tau$. Если при измельчении
сетки результат не стабилизируется, выводы физического характера делать
преждевременно.

\subsection{Верификация, валидация и достоверность}

Следует различать:
\begin{itemize}
    \item \textbf{верификацию} --- правильно ли решаются уравнения выбранной
    математической модели;
    \item \textbf{валидацию} --- достаточно ли хорошо сама модель описывает
    реальный объект в заданной области применимости.
\end{itemize}

Для верификации используют аналитические и эталонные решения, сеточную
сходимость, балансы и инварианты. Для валидации необходимы независимые
данные реального объекта. Поэтому высокая точность вычислений может
означать лишь точное решение неадекватной модели.

\subsection{Преимущества и ограничения}

Вычислительный эксперимент особенно полезен, если натурный эксперимент
дорог, опасен, длителен или принципиально невозможен. Он хорошо
воспроизводим и позволяет быстро исследовать широкую область параметров.
Однако достоверность результата ограничена качеством математической
модели, исходных данных, численного метода и программной реализации.

\subsection{Что важно сказать на экзамене}

Ключевая схема --- \textbf{модель--алгоритм--программа--серия расчётов--анализ}.
Нужно объяснить, что вычислительный эксперимент обязательно включает
контроль сеточной и временной сходимости, исследование параметров и
последующую верификацию/валидацию. Это не просто запуск одной программы.

\newpage
